\documentclass[11pt]{article}
\usepackage[margin=1.15in]{geometry}
\usepackage{amsmath,amssymb,booktabs,tabularx}
\usepackage{microtype,xcolor,mdframed}
\usepackage[colorlinks=true,linkcolor=blue!50!black,urlcolor=blue!50!black]{hyperref}
\newmdenv[backgroundcolor=blue!4,linecolor=blue!35,linewidth=0.8pt,
  skipabove=9pt,skipbelow=9pt,innerleftmargin=10pt,innerrightmargin=10pt]{keyidea}
\newcommand{\code}[1]{{\small\ttfamily #1}}
\setlength{\emergencystretch}{1.5em}
\title{Infinitely Many Collatz Diaries Ruled Out:\\
\Large A Friendly Guide to Short Rotation Windows}
\author{Companion to ``Short Rotation Windows and an Infinite Family\\of Collatz Itinerary Exclusions''\thanks{Written with Codex (OpenAI). The companion technical paper and Lean sources give the exact statements and proofs. No claim of literature priority is made.}}
\date{September 4, 2026}
\begin{document}
\maketitle

\section{What has been proved}
Start with a whole number. Halve it when it is even; when it is odd,
triple it, add one, and halve the result. Record a $1$ at each odd number
and a $0$ at each even number. This infinite binary diary describes its
shortcut Collatz orbit.

There is another way to make a binary diary. Walk along the number line
in steps of a fixed length $\alpha$ between zero and one. Write $1$ when
the step crosses into the next unit interval and $0$ otherwise. The
starting position $\rho$ can be any real number. These are
\emph{mechanical words}; irrational step lengths give Sturmian words.

The previous result ruled out every starting position for one step length,
$1/\sqrt2$. We have now proved an infinite family. For each integer
$m\ge6$, set
\[
t_m=\frac{\sqrt{m^2+4}-m}{2},\qquad \alpha_m=\frac1{1+t_m}.
\]
\begin{keyidea}
\textbf{The new theorem.} No whole-number Collatz orbit has a mechanical
diary of step length $\alpha_m$, for any $m\ge6$ or starting position
$\rho$. An orbit cannot acquire such a diary after a finite beginning,
either. Lean has checked the theorem for all these parameters at once.
\end{keyidea}
Some step lengths, rounded for illustration, are
\begin{center}
\begin{tabular}{rrrrrr}
\toprule
$m$&6&7&8&10&100\\
$\alpha_m$&0.860380&0.877151&0.890388&0.909902&0.990100\\
\bottomrule
\end{tabular}
\end{center}
Lean also proves that they are distinct, irrational, and get arbitrarily
close to one. This is a family of exclusions, not a proof that every
Collatz orbit reaches one.

\newpage
\section{The short-window trick}
Imagine keeping only the fractional position of each point on the number
line. The positions now lie around a circle of circumference one.
Certain repeated diary blocks are located by visits to a short interval
just before position one:
\[
[1-\delta,\,1).
\]
The proof needs to know that a visit happens soon enough, whatever the
starting position. Earlier work used a longer rational grid to guarantee
this. A simpler argument now gives the shorter window $q+Q$.

Choose two approximate returns of the rotation. After $q$ steps, the
position moves forward by a small amount $\delta$ modulo whole turns.
After $Q$ steps, it moves backward by $\varepsilon$, where
$0<\varepsilon\le\delta$.

Look at the first $q+Q$ positions. Suppose that none is in the target
interval. Pick the largest fractional position among them, and let its
time index be $n$. There are two possibilities:
\begin{center}
\begin{tabular}{lll}
\toprule
Index of the largest point&Replacement index&Forward movement\\
\midrule
$n<Q$&$n+q$&$\delta$\\
$n\ge Q$&$n-Q$&$\varepsilon$\\
\bottomrule
\end{tabular}
\end{center}
Both replacement indices are still among the first $q+Q$ indices.
Since the chosen point is below $1-\delta$, neither movement reaches one
or wraps around the circle. Both produce a larger fractional position.
That contradicts the choice of the largest one.

\begin{keyidea}
\textbf{The conclusion.} Every block of $q+Q$ consecutive rotation
positions contains a target visit. The proof works for every starting
position and handles the interval endpoints exactly.
\end{keyidea}
This is all the rotation coverage the Collatz argument needs. We do not
need to classify every possible gap on the circle or assume a
three-distance theorem. A separate integer identity guarantees that
successive visits cannot be too close together. Together the two facts
locate long repeated diary stretches early enough to be useful.

\newpage
\section{Why one recurrence gives infinitely many cases}
For a fixed integer $m\ge6$, begin with four integers
$(q,Q,p,P)=(1,1,0,1)$. Repeat the rules
\[
\begin{aligned}
q'&=q+mQ,&p'&=p+mP,\\
Q'&=mq'+Q,&P'&=mp'+P.
\end{aligned}
\]
They supply the two approximate returns used on the previous page.
The crucial small identity is
\[
t_m^2+mt_m=1.
\]
It says that the next lower error is $t_m^2$ times the previous lower
error; the upper error is always $t_m$ times the lower one. Thus both
errors stay positive, become smaller, and have the required ordering.
The recurrence also preserves the exact integer identity $Pq-pQ=1$,
which provides separation between visits.

For $m=6$, the first useful levels are
\begin{center}
\begin{tabular}{rrrr}
\toprule
Level&Lower return $q$&Separation $Q$&Window $q+Q$\\
\midrule
1&7&43&50\\
2&265&1633&1898\\
3&10063&62011&72074\\
\bottomrule
\end{tabular}
\end{center}
The table is an illustration. Lean proves the recurrence properties at
every level and for every integer parameter $m\ge6$.

\section{The numerical margin is exact}
Matching long parity diaries forces an integer divisibility condition:
two distinct starts sharing their first $L$ parity letters differ by a
nonzero multiple of $2^L$. Their difference is therefore at least $2^L$.
Long repetitions inside one orbit turn this into a lower bound on orbit
size. The exact Collatz formula, including the added-one corrections,
gives an upper bound for the same values.

The short window makes these bounds incompatible at a sufficiently
large recurrence level. One exact comparison does most of the work:
\[
3^8=6561<8192=2^{13}.
\]
After the other factors are bounded, the required estimate is
\[
24\bigl(3N+3(m+2)q+1\bigr)
\left(\frac{6561}{8192}\right)^q<1.
\]
For each fixed starting number $N$ and parameter $m$, the factor outside
the power grows only linearly in $q$. The power shrinks geometrically.
Since the recurrence produces arbitrarily large $q$, the estimate must
eventually hold. There is no finite search limit in this argument.

\newpage
\section{What Lean checked}
The formal proof covers the short-window argument, the determinant
identity, the integer recurrence, the errors and denominator ratios,
the orbit-size comparison, and the final limiting argument. It quantifies
over every starting number, real intercept, finite orbit prefix, and
integer parameter in the stated range.

Lean also checks that the slopes are irrational and strictly increasing,
with the explicit bound
\[
0<1-\alpha_m<\frac1{m+1}.
\]
For example, taking larger and larger $m$ makes the limiting proportion
of odd letters approach $100\%$. Even these highly odd-biased diaries
cannot occur when they have this exact mechanical structure.

The expanded theorem audit checks 26 selected declarations. Their axiom
dependencies are within Lean's standard foundations. This check is not
an independent replay by another kernel or a fresh rebuild of every
external library. The result uses no unproved mathematical axiom and no
unfinished proof placeholder.

\section{What still separates this from Collatz}
\begin{keyidea}
\textbf{The remaining gap.} We can now cross infinitely many slopes,
each with every starting position, off the list of possible diaries.
We have not shown that every hypothetical counterexample must use a
diary on that list.
\end{keyidea}
An infinite set of excluded slopes is not an excluded interval. Getting
arbitrarily close to slope one does not rule out neighboring slopes by
continuity. The theorem does not say that every sufficiently odd-biased
diary is impossible. Nonmechanical diaries, critical-boundary divergence,
and the general nontrivial-cycle problem still need separate arguments.

The next research target is to allow controlled errors in repeated diary
blocks. Earlier work already gives an exact weighted cost for swapping
adjacent letters. A useful theorem must control those weights, not just
count the swaps. It must eventually connect to an independently proved
restriction on the diary of an actual integer counterexample.

\subsection*{Where to read the proof}
The technical companion is \code{paper/metallic.pdf}, with editable
source \code{paper/metallic.tex}. The principal theorem is
\code{Collatz.Metallic.no\_\allowbreak{}itinerary}, in
\code{lean/Collatz/SturmianMetallic.lean}. The supporting arguments are in
\code{SturmianWindow.lean} and
\code{SturmianFamily.lean}. The ongoing research ledger is
\code{analysis/RESEARCH\_\allowbreak{}STATUS.md}.
\end{document}
